% Technical draft. LEE: rewrite the abstract, the introduction and the remarks in Section 7 in your own words.
\documentclass[11pt]{article}
\usepackage[margin=2.5cm]{geometry}
\usepackage{amsmath,amssymb,amsthm}
\usepackage[hidelinks]{hyperref}
\newtheorem{theorem}{Theorem}
\newtheorem{lemma}[theorem]{Lemma}
\newtheorem{proposition}[theorem]{Proposition}
\newcommand{\R}{\mathbb{R}}
\newcommand{\inter}{\operatorname{int}}
\newcommand{\conv}{\operatorname{conv}}
\title{Every convex body in $\mathbb{R}^3$ can be covered by $13$ smaller homothetic copies}
\author{Lee Marshall}
\date{November 2026}
\begin{document}
\maketitle

\begin{abstract}
Every convex body in $\R^3$ is covered by $13$ translates of its interior, so Hadwiger's covering number
satisfies $H_3\le 13$. The previous bound was $14$ (Prymak). We keep Prymak's framework and add one ingredient:
a minimal circumscribed parallelotope forces a chord of length at least $0.569$ (per coordinate) in a
face-diagonal direction for every axis. One of the $13$ covering sets uses this chord. The proof is completed by
an exact linear-programming certificate with $7{,}561{,}583$ leaves, checked by two separately written exact
checkers. The reduction from the certificate to the theorem is formalised in Lean~4. We also correct two points
in the $H_3\le 14$ argument: the symmetry-reduced region used there is not a fundamental domain, and its
Lemma~2.4(ii) fails at some box corners.
\end{abstract}

\section{Introduction}
For a convex body $K\subset\R^n$ let $C(K,\inter K)$ be the least number of translates of $\inter K$ covering $K$,
and $H_n=\max_K C(K,\inter K)$. The conjecture of Levi, Hadwiger, Gohberg and Markus is $H_n=2^n$. In
dimension~$3$ the bounds $20$ (Lassak \cite{Las}), $16$ (Papadoperakis \cite{Pap}) and $14$ (Prymak \cite{Pry}) are known; see \cite{BK} for a survey and \cite{PS,ABP} for related work.
Prymak notes that $14$ is the limit of his method: at the central configuration there are $14$ points no two of
which lie in one translate of his covering body. We remove this obstruction.

\begin{theorem}\label{thm:main}
Every convex body $K\subset\R^3$ is covered by $13$ translates of $\inter K$. Consequently $K$ is covered by $13$
sets $x_i+rK$ with $0<r<1$, and $H_3\le 13$.
\end{theorem}

\section{Framework}
After an affine map, $C=[0,1]^3$ is a parallelotope of minimal volume containing $K$. One can choose contact
points $q_{i0}\in K\cap\{x_i=0\}$ with $q_{i1}=q_{i0}+e_i\in K$ (Papadoperakis, Lemma~3(a)). They are the columns
of
\[
A_p=\begin{pmatrix}0&1&p_3&p_3&p_5&p_5\\ p_1&p_1&0&1&p_6&p_6\\ p_2&p_2&p_4&p_4&0&1\end{pmatrix},
\qquad V_p=\{A_pe_1,\dots,A_pe_6\},\qquad O_p=\conv V_p\subseteq K,
\]
for some $p\in[0,1]^6$. Let $E$ be the edge skeleton of $C$.

\begin{lemma}[Papadoperakis, Lemma~4]\label{lem:red}
Suppose $E\cup V_p\subseteq\bigcup_{i\le m}(v_i+X_i)$, where every $X_i\subseteq\inter K$. Then
$C(K,\inter K)\le m$.
\end{lemma}

We prove both lemmas directly (first variation of the box volume for Lemma~3(a); for each $x\in\partial K$ the
cone $x-\mathrm{cone}(K-x)$ meets $E\cup V_p$, which gives illumination directions). Both proofs are formalised.

\section{The long-chord lemma}
\begin{lemma}\label{lem:M}
Let $[0,1]^3$ be a minimal-volume parallelotope containing $K$. For each axis $i$, with $\{j,k\}$ the other two,
there is $s\in\{\pm1\}$ such that $K$ contains a segment of vector $\tau(e_j+se_k)$, where
$\tau=\tfrac12+\tfrac{\sqrt2-1}{6}=0.5690\ldots$.
\end{lemma}
\begin{proof}
Let $D=K-K$, so $\pm e_l\in D\subseteq[-1,1]^3$. Take $i=1$. For $u_1=e_1$, $u_2=(a,1,1)$, $u_3=(b,1,-1)$ we have
$|\det U|=2$, and the parallelotope cut out by the three slabs of $K$ has volume
$h_D(u_1)h_D(u_2)h_D(u_3)/2$. Minimality gives $h_D(a,1,1)\,h_D(b,1,-1)\ge2$ for all $a,b$, so, say,
$\inf_a h_D(a,1,1)\ge\sqrt2$. Let $(0,t,t)\in\partial D$ with outer normal $n=(a',1+d,1-d)$, $h_D(n)=2t$. Then
$2t\ge1+|d|$ and $\sqrt2\le h_D(n)+h_D(0,-d,d)\le 2t+2|d|$, so $2t\ge1+(\sqrt2-1)/3$.
\end{proof}

With $w=e_j+se_k$, $\lambda=7/10$ and $t=569/1000<\tau$: if $[z,z+tw]\subseteq K$ then
$\lambda z+(1-\lambda)\inter O_p+\lambda t[0,w]\subseteq\inter K$. The position $z$ only shifts this set.

\section{Boxes and strict leaves}
For a box $P\subset[0,1]^6$ let $Q_P=\bigcap_{v}O_v$ over the vertices $v$ of $P$. Prymak's Lemma~2.4(i) gives
$Q_P\subseteq O_p$ for $p\in P$, hence $\inter Q_P\subseteq\inter K$. A \emph{plain leaf} is a box with $13$
translation vectors such that $13$ translates of $\inter Q_P$ cover $E$ and the six rectangles
$\{A_pe_f:p\in P\}$, in a fixed combinatorial pattern (eight vertex translates, five face translates, breakpoints
on edges). A \emph{chord leaf} for axis $i$ has, for each sign $s=\pm1$, such a cover in which one translate is
replaced by $T+\mu\lambda t w_s+(1-\lambda)\inter Q_P$, $\mu\in[0,1]$. All memberships are \emph{strict} linear
inequalities against the exact H-representation of $Q_P$. By Lemmas~\ref{lem:red} and~\ref{lem:M}, if $p$ lies in a
leaf then $C(K,\inter K)\le 13$.

Prymak instead uses Lemma~2.4(ii), that some translate of $Q_P$ lies in $\inter O_p$. This fails for
$P=\prod[a_i,a_i+\tfrac1{160}]$, $a=(0,0,\tfrac{63}{160},\tfrac58,0,\tfrac{57}{80})$, at the corner
$p=a+\tfrac1{160}(1,1,0,0,1,0)$: $Q_P$ meets facets of $O_p$ with inner normals
$n_1=(0,80,9040)$, $n_2=(18080,-160,0)$, $n_3=(-480,0,-496)$, $n_4=(3680,-4960,0)$, and
$\tfrac{215760}{113}n_1+\tfrac{104377}{113}n_2+34800\,n_3+n_4=0$. So every translate of $Q_P$ has a point outside
$\inter O_p$. Strict leaf inequalities make Lemma~2.4(ii) unnecessary.

\section{Symmetry}
The cube symmetries act on $p$ by signed coordinate permutations. Let $s_x=p_3+p_5$, $s_y=p_1+p_6$,
$s_z=p_2+p_4$. Prymak's region $D=\{p_1,p_2,p_3\le\tfrac12,\ s_z\le s_y\le s_x\}$ is not a fundamental domain:
$p=(\tfrac{459}{500},\tfrac{1481}{5000},\tfrac{1689}{10000},\tfrac{9529}{10000},\tfrac{9359}{10000},\tfrac{3079}{10000})$
has no image in $D$ or in the union of the start boxes certified in~\cite{Pry}; about $2.7\%$ of random $p$ behave
this way. The reason is that a reflection sends $s_i\mapsto 2-s_i$, so permutations can order $|s_i-1|$ but not
$s_i$. The region $D^*=\{p_1,p_2,p_3\le\tfrac12,\ |s_z-1|\le|s_y-1|\le|s_x-1|\}$ meets every orbit.

We do not certify $D^*$. Symmetries map the $1/10$ grid of $[0,1]^6$ to itself, so it suffices that every orbit
of $1/10$ boxes contains a certified start box. The $39{,}380$ start boxes of Prymak's region miss $570$ orbits;
we certify one box of each. Given $K$, a symmetry $g$ moves its configuration into a certified box, and the leaf is
applied to $gK$.

\section{The certificate and its checks}
\begin{itemize}
\item $39{,}950$ start boxes of side $1/10$, bisected on the longest side to side at least $1/640$;
  $7{,}561{,}583$ leaves. The $999$ chord leaves all lie in the $8$ start boxes around $(\tfrac12)^6$.
\item Solutions are rational with denominator $10^9$, found by floating-point LPs and stored per leaf.
\item Checker 1 (\texttt{certcheck\_strict.py}): exact integer arithmetic, strict inequalities, tiling of every
  start box; all leaves pass. A non-strict variant, cross-checked against a rational reference on $1\%$ of the
  leaves, also passes. A separate script checks orbit coverage of the $10^6$ grid boxes.
\item Checker 2 (\texttt{hcheck.py}): written from the specification only; ignores the stored structures and
  requires the exact intervals of the $13$ translates on each edge to cover it; all leaves pass strictly, all
  start boxes are tiled, and every grid orbit contains a certified box.
\item Lean~4: Lemmas~3(a) and~4 of Papadoperakis, existence of the minimal box, Lemma~\ref{lem:M}, Lemma~2.4(i),
  leaf soundness, the symmetry reduction and the passage to homothets are proved without \texttt{sorry}. The only
  assumptions beyond Lean's standard axioms are the certificate (every certified start box is tiled by leaves
  satisfying the strict leaf predicate) and one \texttt{native\_decide} computation (orbit coverage of the grid).
\end{itemize}

\section{Remarks}
The two corrections apply to the proof of $H_3\le14$ in~\cite{Pry}. The $570$ extension boxes are certified by
plain leaves, which lie inside Prymak's framework, and the strict-containment argument replaces Lemma~2.4(ii) for
every leaf that satisfies its inequalities strictly. We have not re-checked Prymak's certificate under the strict
test. Lemma~\ref{lem:M} is elementary and may be known; we did not find it in the literature.

\section*{Data}
\begin{itemize}
\item Code, checkers, Lean project and certificate hashes:
  \url{https://github.com/LeeMarshall1113/ai-math-results}, folder \texttt{results/hadwiger-h3}; archived as
  doi:10.5281/zenodo.23748157 (all versions).
\item Certificate data ($2.2$\,GB): doi:10.5281/zenodo.XXXXXXXX. SHA-256 of the main archive:\
  \texttt{8929c5b3599230d0d63234d0497d06e0eda12741380b4363666309987bd995bb}.
\end{itemize}

\section*{AI use}
This work was carried out with AI assistance (Claude, Anthropic), directed by the author. The certificate was
checked by two separately written exact checkers, the reduction is formalised in Lean, and the result was attacked
by adversarial AI referees. It has not been peer reviewed. The author is responsible for all content.

\begin{thebibliography}{9}
\bibitem{Pap} I.~Papadoperakis, An estimate for the problem of illumination of the boundary of a convex body in
  $E^3$, \emph{Geom.\ Dedicata} 75 (1999) 275--285, doi:10.1023/A:1005056207406.
\bibitem{Pry} A.~Prymak, A new bound for Hadwiger's covering problem in $\mathbb{E}^3$, \emph{SIAM J.\ Discrete
  Math.} 37 (2023) 17--24, doi:10.1137/22M1490314.
\bibitem{PS} A.~Prymak and V.~Shepelska, On the Hadwiger covering problem in low dimensions, \emph{J.\ Geom.}
  111 (2020) 42, arXiv:1811.08962.
\bibitem{ABP} A.~Arman, A.~Bondarenko and A.~Prymak, On Hadwiger's covering problem in small dimensions,
  \emph{Canad.\ Math.\ Bull.} (2025), arXiv:2404.00547.
\bibitem{Las} M.~Lassak, Covering a three-dimensional convex body by smaller homothetic copies,
  \emph{Beitr\"age Algebra Geom.} 39 (1998) 259--262.
\bibitem{BK} K.~Bezdek and M.~A.~Khan, The geometry of homothetic covering and illumination, in \emph{Discrete
  Geometry and Symmetry}, Springer (2018) 1--30.
\end{thebibliography}
\end{document}
